\chapter{}
\label{chap:3}

\section{Умова}

Оцінити часову складність атаки Куртуа-Майєра на фільтрувальний генератор
гами, побудований на основі ЛРЗ довжини $n$ з функцією ускладнення $f$, якщо
\begin{enumerate}[label=\alph*)]
    \item $f(x_{1}, x_{2}, x_{3}) = x_{1} \oplus x_{2}x_{3}$;
    \item $f(x_{1}, x_{2}, x_{3}) = x_{2} \oplus x_{1}x_{2} \oplus x_{1}x_{3}$.
\end{enumerate}

Виписав спершу трохи теорії для себе для зручності:

Розглянемо фільтрувальний генератор гами, що описується системою рівнянь:
\begin{equation}
    \label{eq:system_of_equation}
    f(xL_i) = \gamma_i, \qquad i = 0, 1, 2, \ldots,
\end{equation}
де $x \in V_{n}$ -- невідомий початковий стан ЛРЗ довжини~$n$, $f \in B_n$ -- обрана функція ускладнення, 
$L_{i}$ -- відома матриця розміру $n \times n$ з двох елементів. Припустимо, що виконується хоча б одна з умов:
\begin{enumerate}[leftmargin=2em]
    \item існує ненульова $g \in \Ann(f)$ така, що $\deg g < \deg f$;
    \item існує ненульова $h \in \Ann(f \oplus 1)$ така, що $\deg h < \deg f$.
\end{enumerate}
Тоді у випадку $\gamma_i = 1$ (+ умова 1) виходить $0 = g(xL_i) \cdot f(xL_i) = g(xL_i)$, а у випадку $\gamma_i = 0$ 
(+ умова 2) маємо $0 = h(xL_i) \cdot (f(xL_i) \oplus 1) = h(xL_i)$.

Позначимо через $d = \min\{\deg g,\; \deg h\}$. Представляємо ліву частину системи~\eqref{eq:system_of_equation} у 
вигляді поліному Жегалкіна (АНФ) та заміняємо кожен моном $x^{\alpha}$ на нову змінну $y_{\alpha}$ (метод введення 
нових змінних). Отримаємо систему лінійних рівнянь від $N_{d}$ нових невідомих, що являє собою суму можливої кількості 
мономів від $1$ до $d$-того степеня (мономи нульового степеня не враховуємо, бо вони лишаються незмінними)
\begin{equation*}
    N_{d} = \binom{n}{1} + \binom{n}{2} + \ldots + \binom{n}{d} = \sum\limits_{i=1}^{d} \binom{n}{i}
\end{equation*}
Складність розв'язання новоутвореної системи методом Гауса складає $O(t \cdot N_d^2)$ двійкових операцій, де 
$t \geq N_d$ -- кількість рівнянь у системі.

Виграш по часу атаки Куртуа--Майєра порівняно з атакою з розв'язанням вхідної системи методом введення нових змінних 
є величиною порядку $(N_{D} / N_{d})^3$, де $D = \deg f$ -- відома (в межах задачі) величина. Число 
$d = \AI(f)$ -- алгебраїчна імунність функції $f$. Тому щоб обчислити складність, необхідно обчислити $\AI(f)$.

\section{Розв'язання}

\subsection*{Пункт а)}
$f(x_1, x_2, x_3) = x_1 \oplus x_2 x_3$, $D = \deg f = 2$.

\noindent Множина нулів та одиниць для даної функції має вигляд:
\begin{center}
    \begin{tabular}{c|cccccccc}
        $(x_1, x_2, x_3)$ & $(0,0,0)$ & $(0,0,1)$ & $(0,1,0)$ & $(0,1,1)$
                          & $(1,0,0)$ & $(1,0,1)$ & $(1,1,0)$ & $(1,1,1)$                         \\
        \hline
        $f$               & $0$       & $0$       & $0$       & $1$       & $1$ & $1$ & $1$ & $0$
    \end{tabular}
\end{center}
\begin{equation*}
    V (I) = \{(0,0,0),\; (0,0,1),\; (0,1,0),\; (1,1,1)\}, \quad
    V_{n} \setminus V(I) = \{(0,1,1),\; (1,0,0),\; (1,0,1),\; (1,1,0)\}.
\end{equation*}
Щоб обчислити алгебраїчну імунність, треба порахувати дві величини:

Для початку обчислюємо $\min \deg \Ann(f)$. Ми доводили, що $\AI (f) \leq \lceil \frac{n}{2}\rceil$, в нашому випадку 
$\AI (f) \leq 2$, а тому перевіряємо, чи існує ненульова функція меншого степеня (в даному випадку $1$) в $\Ann(f)$. 
Нехай $g = a_0 \oplus a_1 x_1 \oplus a_2 x_2 \oplus a_3 x_3 \in \Ann(f)$, де $a_{i}$ -- невідомі коефіцієнти. $g$ 
має обертатися в нуль на ненульових виходах $V_{n} \setminus V(I)$:
\begin{equation*}
    \begin{cases}
        g(0,1,1) = a_0 \oplus a_2 \oplus a_3 = 0, \\
        g(1,0,0) = a_0 \oplus a_1 = 0, \\
        g(1,0,1) = a_0 \oplus a_1 \oplus a_3 = 0 \;\Rightarrow\; a_3 = 0, \\
        g(1,1,0) = a_0 \oplus a_1 \oplus a_2 = 0 \;\Rightarrow\; a_2 = 0.
    \end{cases}
\end{equation*}
З ксору третього та другого рівнянь: $a_{3} = 0$. З четвертого та другого: $a_{2} = 0$. Тоді з першого $a_0 = 0$, 
і з другого: $a_1 = 0$. Отже маємо єдиний розв'язок для $g$, який є тривіальним, тобто ненульових ануляторів степеня 
$1$ в $\Ann(f)$ не існує. А оскільки $f \oplus 1 \in \Ann(f)$ і $\deg(f \oplus 1) = 2$, маємо $\min \deg \Ann(f) = 2$.

Далі обчислимо $\min \deg \Ann(f \oplus 1)$. Аналогічно, нехай $h = a_0 \oplus a_1 x_1 \oplus a_2 x_2 \oplus a_3 x_3 \in \Ann(f \oplus 1)$, 
тобто $h$ обертається в нуль на $V(I)$:
\begin{equation*}
    \begin{cases}
        h(0,0,0) = a_0 = 0, \\
        h(0,0,1) = a_3 = 0, \\
        h(0,1,0) = a_2 = 0, \\
        h(1,1,1) = a_1 \oplus a_2 \oplus a_3 = a_1 = 0.
    \end{cases}
\end{equation*}
Єдиний розв'язок є тривіальним, тобто ненульових ануляторів степеня $1$ в $\Ann(f \oplus 1)$ також не існує. 
Оскільки $f \in \Ann(f \oplus 1)$ і $\deg f = 2$, маємо $\min \deg \Ann(f \oplus 1) = 2$.

\noindent Отже алгебраїчна імунність буде мінімумом з двох значень:
\begin{equation*}
    \AI(f) = \min\big\{\min \deg \Ann(f),\; \min \deg \Ann(f \oplus 1)\big\} = \min\{2,\; 2\} = 2.
\end{equation*}
Тоді для оцінки часової складності, при $d = \AI(f) = 2$, нам треба порахувати:
\begin{equation*}
    N_{d} = N_{2} = \binom{n}{1} + \binom{n}{2} = n + \frac{n(n-1)}{2} = \frac{n(n+1)}{2}.
\end{equation*}
Необхідна кількість бітів гами: $t \geq N_2 = \dfrac{n(n+1)}{2}$. А часова складність атаки:
\begin{equation*}
    T \approx O \left(t \cdot N_2^2\right) = O \left(N_2^3\right) = O\!\left(\frac{n^3(n+1)^3}{8}\right) \approx O(n^6).
\end{equation*}
Оскільки отримали $d = D = 2$, то атака Куртуа--Майєра не дає виграшу в часовій складності порівняно з
прямим розв'язанням вхідної системи методом введення нових змінних $(\left(\dfrac{N_D}{N_d}\right)^{3} = 1)$

\subsection*{Пункт б)}

$f(x_1, x_2, x_3) = x_2 \oplus x_1 x_2 \oplus x_1 x_3$, $D = \deg f = 2$.

$f(x_1, x_2, x_3) = x_1 \oplus x_2 x_3$, $D = \deg f = 2$.

\noindent Множина нулів та одиниць для даної функції має вигляд:
\begin{center}
    \begin{tabular}{c|cccccccc}
        $(x_1, x_2, x_3)$ & $(0,0,0)$ & $(0,0,1)$ & $(0,1,0)$ & $(0,1,1)$
                          & $(1,0,0)$ & $(1,0,1)$ & $(1,1,0)$ & $(1,1,1)$                         \\
        \hline
        $f$               & $0$       & $0$       & $1$       & $1$       & $0$ & $1$ & $0$ & $1$
    \end{tabular}
\end{center}
\begin{equation*}
    V (I) = \{(0,0,0),\; (0,0,1),\; (1,0,0),\; (1,1,0)\}, \quad
    V_{n} \setminus V(I) = \{(0,1,0),\; (0,1,1),\; (1,0,1),\; (1,1,1)\}.
\end{equation*}
Алгебраїчну імунність, обчислюємо так само, знаходячи дві величини:

Обчислюємо $\min \deg \Ann(f)$. Перевіримо, чи існує ненульова функція степеня~$1$ в $\Ann(f)$. Нехай 
$g = a_0 \oplus a_1 x_1 \oplus a_2 x_2 \oplus a_3 x_3 \in \Ann(f)$, тобто $g$ має обертається в нуль на $V_{n} \setminus V(I)$:
\begin{equation*}
    \begin{cases}
        g(0,1,0) = a_0 \oplus a_2 = 0,                                    \\
        g(0,1,1) = a_0 \oplus a_2 \oplus a_3 = 0 \;\Rightarrow\; a_3 = 0, \\
        g(1,0,1) = a_0 \oplus a_1 \oplus a_3 = a_0 \oplus a_1 = 0,        \\
        g(1,1,1) = a_0 \oplus a_1 \oplus a_2 \oplus a_3 = a_0 \oplus a_1 \oplus a_2 = 0.
    \end{cases}
\end{equation*}
Проксоривши 1 і 2, маємо $a_3 = 0$, друге і четверте -- $a_1 = 0$. З третього тоді: $a_0 = 0$, і з першого 
$a_0 = 0$. Тобто єдиний розв'язок --- тривіальний і ненульових ануляторів степеня $1$ в $\Ann(f)$ не існує. 
Оскільки $f \oplus 1 \in \Ann(f)$ і $\deg(f \oplus 1) = 2$, маємо $\min \deg \Ann(f) = 2$.

Далі рахуємо $\min \deg \Ann(f \oplus 1)$. Нехай $h = a_0 \oplus a_1 x_1 \oplus a_2 x_2 \oplus a_3 x_3 \in \Ann(f \oplus 1)$, 
тобто $h$ обертається в нуль на $V(I)$:
\begin{equation*}
    \begin{cases}
            h(0,0,0) = a_0 = 0, \\
            h(0,0,1) = a_3 = 0, \\
            h(1,0,0) = a_1 = 0, \\
            h(1,1,0) = a_1 \oplus a_2 = a_2 = 0.
    \end{cases}
\end{equation*}
Тут одразу отримаємо єдиний розв'язок, що є тривіальним, тобто ненульових ануляторів степеня $1$ в $\Ann(f \oplus 1)$ 
також не існує. Таким чином, оскільки $f \in \Ann(f \oplus 1)$ і $\deg f = 2$, маємо $\min \deg \Ann(f \oplus 1) = 2$.

\noindent Алгебраїчна імунність теж виходить рівна двом:
\begin{equation*}
    \AI(f) = \min\big\{\min \deg \Ann(f),\; \min \deg \Ann(f \oplus 1)\big\} = \min\{2,\; 2\} = 2.
\end{equation*}
Для часової складності $d = \AI(f) = 2$ обчислюємо:
\begin{equation*}
    N_d = N_2 = \binom{n}{1} + \binom{n}{2} = \frac{n(n+1)}{2}.
\end{equation*}
Необхідна кількість бітів гами складає: $t \geq N_2 = \dfrac{n(n+1)}{2}$. А часова складність
\begin{equation*}
    T \approx O \left(t \cdot N_2^2\right) = O \left(N_2^3\right) \approx O(n^6).
\end{equation*}
Виграшу, порівняно зі звичайною атакою теж немає, $(N_{D}/N_{d})^{3} = 1$, оскільки $\AI(f) = \deg f = 2$.

В обох пунктах $\AI(f) = D = 2$ алгебраїчна імунність збігається зі степенем функції ускладнення, тому виграшу 
саме від пошуку ануляторів меншого степеня немає. Тим не менш, складність $O(n^6)$ є суттєво меншою за складність 
повного перебору $O(2^n)$, яка є експоненціальною.
